% zarrabi/thesis-template-inspired Persian starter. Compile with XeLaTeX.
\documentclass[12pt,a4paper,oneside]{book}
\usepackage[a4paper,margin=30mm]{geometry}
\usepackage{polyglossia}
\setmainlanguage{farsi}
\setotherlanguage{english}
% FreeSerif covers Persian digits and Arabic-script glyphs in the pinned image;
% DejaVu Serif does not cover the full Persian range.
\newfontfamily\persianfont{FreeSerif}
\newfontfamily\englishfont{FreeSerif}
\usepackage{booktabs}
\usepackage{amsmath}
\usepackage[hidelinks]{hyperref}
\title{قالب پایان‌نامه فارسی\\Persian Thesis Starter}
\author{Alex Morgan}
\date{۱۴۰۵}
\begin{document}
\frontmatter
\maketitle
\chapter*{چکیده}
\addcontentsline{toc}{chapter}{چکیده}
این یک قالب آغازین مستقل برای پایان‌نامه‌های فارسی است. عنوان، مشخصات نویسنده،
فصل‌ها، فهرست مطالب و پیوست‌ها را مطابق آیین‌نامه دانشگاه خود تنظیم کنید.
\chapter*{Abstract}
\addcontentsline{toc}{chapter}{Abstract}
\begin{english}
This starter demonstrates a Persian thesis workflow with right-to-left front
matter and ordinary Latin technical content.
\end{english}
\tableofcontents
\mainmatter
\chapter{مقدمه}
در این فصل مسئله پژوهش، پیشینه و سهم علمی کار را معرفی کنید.
\chapter{روش پژوهش}
\begin{english}
An illustrative model is
\begin{equation}
  y = \alpha + \beta x + \varepsilon.
\end{equation}
\begin{table}[ht]
  \centering
  \caption{Illustrative measurements.}
  \begin{tabular}{@{}lrr@{}}
    \toprule Case & Estimate & Error \\
    \midrule A & 1.24 & 0.08 \\
    B & 1.51 & 0.05 \\
    \bottomrule
  \end{tabular}
\end{table}
\end{english}
\chapter{نتیجه‌گیری}
نتایج، محدودیت‌ها و مسیرهای ادامه پژوهش را در این فصل خلاصه کنید.
\appendix
\chapter{پیوست}
جداول و جزئیات تکمیلی را در پیوست قرار دهید.
\end{document}
